\section{\HMTLang{Composition de l’action, de la fréquence propre et de la réciprocité radiale}{Composition de l’action, de la fréquence propre et de la réciprocité radiale}
}
\label{delta22:vi:composicion}
\HMTLang{L’opérateur de masse précédent reçoit la signature, le caractère et les tours de l’histoire enrichie. La section d’action fixe son échelle absolue et participe également au contre-angle et au correcteur de retour. La composition suivante explicite cette conjonction : une même coordonnée structurelle de route détermine la masse, la fréquence propre, deux lectures radiales réciproques et, dans la réalisation canonique spécifiée, la pondération thermique.}{L’opérateur de masse précédent reçoit la signature, le caractère et les tours de l’histoire enrichie. La section d’action fixe son échelle absolue et participe également au contre-angle et à la correction de retour. La composition suivante explicite cette conjonction : une coordonnée structurelle de route détermine la masse, la fréquence propre, deux lectures radiales réciproques et, dans la réalisation canonique spécifiée, la pondération thermique.}
\HMTLang{L’ordre APP--TRIT--TPK, état enrichi et structure discrète conjointe du continu est conservé ; les coordonnées régionales et \(\alpha\) sont des sorties précédemment engendrées. La sélection de \(K\), l’incidence et l’orientation restent dans leurs antécédents. Les opérations de ce paragraphe sont des compositions ultérieures des lecteurs déjà construits.}{L’ordre APP--TRIT--TPK, état enrichi et structure discrète conjointe du continu est conservé ; les coordonnées régionales et \(\alpha\) sont des sorties précédemment engendrées. La sélection de \(K\), l’incidence et l’orientation restent dans leurs antécédents. Les opérations de cette section sont des compositions ultérieures des lecteurs déjà construits.}


\subsection{\HMTLang{Restitution radiale de la section d’action}{Récupération radiale de la section d’action}
}
\HMTLang{Avec la paire angulaire dans une même unité, nous écrivons}{Avec la paire angulaire dans la même unité, écrivons}

\[
 A=1000\alpha,\quad C^*=2(\eta_{\rm ret}+\alpha),\quad
 \xi=C^*/A,\quad
 R_\star=\left(\frac{1-\xi}{1+\xi}\right)^{1/4},\quad q=R_\star^4.
\]
\HMTLang{Sur la branche positive \(0<\eta_{\rm ret}<499\alpha\), on a \(0<q<499/501\). L’inversion explicite est}{Sur la branche positive \(0<\eta_{\rm ret}<499\alpha\), on a \(0<q<499/501\). L’inverse explicite est}

\begin{equation}
 \eta_{\rm ret}=\alpha\frac{499-501q}{1+q},\qquad
 \hbar_{\rm ret}=10^{-34}\mathcal U_S\eta_{\rm ret},\qquad
 R_{\rm act}=\frac{H_5(1+q)}{\alpha(499-501q)}.
 \label{delta22:vi:accion}
\end{equation}
\begin{proof}
\HMTLang{L’identité \(q=(1-\xi)/(1+\xi)\) donne \(\xi=(1-q)/(1+q)\). La substituer dans \(\eta_{\rm ret}=\alpha(500\xi-1)\) prouve la première formule ; les deux autres utilisent les définitions de la section d’action et du correcteur \(R_{\rm act}=H_5/\eta_{\rm ret}\). Pour les deux orientations, soit \(\varepsilon=\pm1\) et conservons \(q_\varepsilon=(1-\varepsilon\xi)/(1+\varepsilon\xi)\). Alors}{L’identité \(q=(1-\xi)/(1+\xi)\) donne \(\xi=(1-q)/(1+q)\). La substitution dans \(\eta_{\rm ret}=\alpha(500\xi-1)\) prouve la première formule ; les deux autres utilisent les définitions de la section d’action et de la correction \(R_{\rm act}=H_5/\eta_{\rm ret}\). Pour les deux orientations, soit \(\varepsilon=\pm1\) et conservons \(q_\varepsilon=(1-\varepsilon\xi)/(1+\varepsilon\xi)\). Alors}

\[
 \eta_{\rm ret}=\alpha\left[500\varepsilon
 \frac{1-q_\varepsilon}{1+q_\varepsilon}-1\right].
\]
\HMTLang{La transformation \((\varepsilon,q_\varepsilon)\mapsto(-\varepsilon,q_\varepsilon^{-1})\) conserve l’expression et, avec elle, \(\hbar_{\rm ret}\) et \(R_{\rm act}\).}{La transformation \((\varepsilon,q_\varepsilon)\mapsto(-\varepsilon,q_\varepsilon^{-1})\) conserve l’expression et donc \(\hbar_{\rm ret}\) et \(R_{\rm act}\).}

\end{proof}
\HMTLang{Cette restitution conserve la marque d’orientation, \(\alpha\), \(H_5\) et la base dimensionnelle. La signature et les tours continuent de provenir du registre de route.}{Cette récupération conserve la marque d’orientation, \(\alpha\), \(H_5\) et la base dimensionnelle. La signature et les tours continuent de provenir du registre de route.}


\subsection{\HMTLang{Section longitudinale, horloge induite et cotransformation électronique}{Section longitudinale, horloge induite et cotransformation électronique}
}
\label{delta22:vi:longitud}
\HMTLang{Soient \(L_\star>0\) la section longitudinale de référence et \(c_{\rm int}>0\) la vitesse constitutive dans la coordinatisation choisie. Nous écrivons \(\alpha_{\rm HMT}\) et \(\pi_{\rm HMT}\) pour les coordonnées déjà engendrées par l’état enrichi ; ce ne sont pas des données métrologiques qui sélectionnent cette construction. Dans les formules précédemment établies qui suivent, \(\alpha:=\alpha_{\rm HMT}\) et \(\pi:=\pi_{\rm HMT}\). L’incidence marquée contient}{Soient \(L_\star>0\) la section longitudinale de référence et \(c_{\rm int}>0\) la vitesse constitutive dans la carte choisie. Nous écrivons \(\alpha_{\rm HMT}\) et \(\pi_{\rm HMT}\) pour les coordonnées déjà engendrées par l’état enrichi ; ce ne sont pas des données métrologiques sélectionnant cette construction. Dans les formules précédemment établies ci-dessous, \(\alpha:=\alpha_{\rm HMT}\) et \(\pi:=\pi_{\rm HMT}\). L’incidence marquée contient}

\begin{equation}
 N=12\cdot4\cdot120=5760,
 \qquad r_N=\frac{N-1}{4N}=\frac{5759}{23040},
 \qquad
 L_\alpha=\alpha_{\rm HMT}^{16}r_NL_\star .
 \label{delta22:vi:Lalpha}
\end{equation}
\HMTLang{Le coefficient \(r_N\) est l’espérance diagonale du retour résiduel de Hadamard, \(\Delta(CC^*)=r_NI\). Le canal complémentaire \(1/(4N)\) est conservé comme complément, mais n’est pas ajouté à l’observable longitudinale. Par conséquent, \(L_\alpha\) est construit avant l’horloge et avant toute lecture gravitationnelle.}{Le coefficient \(r_N\) est l’espérance diagonale du retour résiduel de Hadamard, \(\Delta(CC^*)=r_NI\). Le canal complémentaire \(1/(4N)\) est conservé comme complément, mais n’est pas ajouté à l’observable longitudinale. Ainsi, \(L_\alpha\) est construit avant l’horloge et avant toute lecture gravitationnelle.}


\HMTLang{L’horloge de \(108\) pas est la réalisation circulaire ultérieure de cette longueur :}{L’horloge à \(108\) pas est la réalisation circulaire ultérieure de cette longueur :}

\begin{equation}
 t_0=\frac{\pi_{\rm HMT}L_\alpha}{54c_{\rm int}},\qquad
 \omega_0=\frac{2\pi_{\rm HMT}}{108t_0}=\frac{c_{\rm int}}{L_\alpha},
 \qquad R_{108}=\frac{c_{\rm int}}{\omega_0}=L_\alpha .
 \label{delta22:vi:reloj-longitudinal}
\end{equation}
\HMTLang{Une section d’action retournée \(\hbar_{\rm ret}>0\) étant fixée, on définit l’énergie et la masse de l’horloge, et alors seulement le facteur électronique :}{Pour une section d’action retournée fixée \(\hbar_{\rm ret}>0\), l’énergie et la masse de l’horloge sont définies d’abord, et alors seulement le facteur électronique :}

\begin{equation}
 Q_L=\frac{\hbar_{\rm ret}c_{\rm int}}{L_\alpha},\qquad
 E_{\rm clk}=Q_L,\qquad
 m_{\rm clk}=\frac{Q_L}{c_{\rm int}^2},\qquad
 E_e^{(0)}=\varepsilon_e^{(0)}U_E,\qquad
 b_e(Q_L)=\frac{E_e^{(0)}}{Q_L}.
 \label{delta22:vi:escala-longitudinal}
\end{equation}
\HMTLang{Dans ce qui suit, \(b_e\) abrège la fonction évaluée \(b_e(Q_L)\).}{Ci-après, \(b_e\) abrège la fonction évaluée \(b_e(Q_L)\).}
\HMTLang{On conserve l’un des lecteurs de retour \(r\in\{D,\Omega\}\). Le caractère raffiné est celui construit dans le paragraphe précédent, avec convergence et tours relatives à l’électron conservées.}{L’un des lecteurs de retour \(r\in\{D,\Omega\}\) est conservé. Le caractère raffiné est celui construit dans la section précédente, en conservant la convergence et les tours relatives à l’électron.}
\HMTLang{La coordonnée électronique basale et son exposant, dérivés précédemment, sont}{La coordonnée électronique basale et son exposant, dérivés ci-dessus, sont}

\[
 d_4=\frac{\pi-e\log\pi}{270},\quad
 \Delta_4=d_4\left(1+\frac\pi{729}\right),\quad
 \varepsilon_e^{(0)}=\frac{\sqrt3}{4}
 \left[e^{\varphi/\pi^2}-22\alpha^3\right](1+15\Delta_4),\quad
 \kappa_e=80-54\alpha+6\Delta_4.
\]
\HMTLang{Pour une route à réalisation scalaire positive,}{Pour une route à réalisation scalaire positive,}

\begin{equation}
 \Xi_\Gamma=b_e\,
 \chi_{\rm ref}(\sigma_\Gamma-\sigma_e,\eta_\Gamma-\eta_e)
 R_{\rm act}^{\kappa_\Gamma^r},\qquad
 m_\Gamma^\beta=m_{\rm clk}\Xi_\Gamma.
 \label{delta22:vi:xi}
\end{equation}
\HMTLang{Cette formule emploie les tours absolues avant de prendre leur différence ; lorsque le registre fournit déjà les tours relatives, cette différence n’est prise qu’une seule fois. Le terme de hauteur \(120\) conserve également son inclusion unique dans le caractère. Dans l’état central,}{Cette formule utilise les tours absolues avant de prendre leur différence ; lorsque le registre fournit déjà les tours relatives, cette différence n’est prise qu’une seule fois. Le terme de hauteur \(120\) conserve de même son inclusion unique dans le caractère. À l’état central,}

\[
 m_e^\beta=\frac{\varepsilon_e^{(0)}U_E}{c_{\rm int}^2}
 R_{\rm act}^{\kappa_e}.
\]
\HMTLang{L’égalité se vérifie en simplifiant \(Q_L\) entre \(m_{\rm clk}\) et \(b_e(Q_L)\). Ainsi, la normalisation électronique complète est conservée. L’action intervient dans l’échelle absolue, dans \(R_{\rm act}\) et dans le contre-angle qui détermine les canaux constitutifs. Aucune correction décimale dépendant de l’espèce n’apparaît : la décade de la section d’action reste exclusivement dans la branche absolue action--horloge--masse.}{L’égalité découle de la simplification de \(Q_L\) entre \(m_{\rm clk}\) et \(b_e(Q_L)\). Cela conserve la normalisation électronique complète. L’action intervient dans l’échelle absolue, \(R_{\rm act}\), et le contre-angle déterminant les canaux constitutifs. Aucune correction décimale dépendant de l’espèce n’apparaît : la décade de la section d’action reste exclusivement dans la branche absolue action--horloge--masse.}


\HMTLang{La covariance de coordinatisation empêche de figer \(b_e\). Avec \(c_{\rm int}\), \(E_e^{(0)}\) et toutes les données adimensionnelles de la route fixés ---caractère raffiné, exposants et \(R_{\rm act}\)---, la cotransformation des sections longitudinale et d’action s’écrit}{La covariance de carte empêche de figer \(b_e\). Avec \(c_{\rm int}\), \(E_e^{(0)}\) et toutes les données adimensionnelles de la route fixés---le caractère raffiné, les exposants et \(R_{\rm act}\)---, la cotransformation des sections longitudinale et d’action s’écrit}

\[
 L_\star\longmapsto\lambda L_\star,\qquad
 \hbar_{\rm ret}\longmapsto\mu\hbar_{\rm ret},\qquad \lambda,\mu>0,
\]
\HMTLang{alors}{then}
\begin{equation}
 L_\alpha\longmapsto\lambda L_\alpha,\qquad
 Q_L,m_{\rm clk}\longmapsto\frac{\mu}{\lambda}(Q_L,m_{\rm clk}),
 \qquad b_e\longmapsto\frac{\lambda}{\mu}b_e,\qquad
 m_{\rm clk}b_e=\frac{E_e^{(0)}}{c_{\rm int}^2}.
 \label{delta22:vi:be-cotransformacion}
\end{equation}
\HMTLang{Par conséquent, la variation des sections cotransforme le facteur électronique et n’altère ni les masses déjà construites ni leurs tableaux.}{Par conséquent, faire varier les sections cotransforme le facteur électronique et n’altère ni les masses déjà construites ni leurs tableaux.}


\HMTLang{Dans les identités suivantes, on abrège \(m_\Gamma=m_\Gamma^\beta\), en conservant la normalisation électronique complète.}{Dans les identités suivantes, \(m_\Gamma=m_\Gamma^\beta\) est utilisé comme abréviation, en conservant la normalisation électronique complète.}

\begin{proposition}[\HMTLang{Fréquence propre et longueur réduite}{Fréquence propre et longueur d’onde réduite}
]
\HMTLang{Avec \(E_\Gamma=m_\Gamma c_{\rm int}^2\), \(\Omega_\Gamma=E_\Gamma/\hbar_{\rm ret}\), \(\nu_\Gamma=\Omega_\Gamma/(2\pi_{\rm HMT})\) et \(R_{108}=L_\alpha\), on a}{Avec \(E_\Gamma=m_\Gamma c_{\rm int}^2\), \(\Omega_\Gamma=E_\Gamma/\hbar_{\rm ret}\), \(\nu_\Gamma=\Omega_\Gamma/(2\pi_{\rm HMT})\) et \(R_{108}=L_\alpha\), on a}

\[
 E_\Gamma=E_{\rm clk}\Xi_\Gamma,\quad
 \Omega_\Gamma=\omega_0\Xi_\Gamma,\quad
 \nu_\Gamma=\frac{\Xi_\Gamma}{108t_0},\quad
 \bar\lambda_\Gamma=\frac{\hbar_{\rm ret}}{m_\Gamma c_{\rm int}}
 =\frac{R_{108}}{\Xi_\Gamma}.
\]
\end{proposition}
\begin{proof}
\HMTLang{Substituer \eqref{delta22:vi:xi} et les définitions de \(m_{\rm clk}\) et de \(E_{\rm clk}\) produit les quatre identités. Pour deux routes positives, le quotient des fréquences est le quotient des masses ; la propriété multiplicative du caractère conserve la signature et les tours relatives, avec \(R_{\rm act}^{\kappa_Y^r-\kappa_X^r}\).}{Substituer \eqref{delta22:vi:xi} et les définitions de \(m_{\rm clk}\) et de \(E_{\rm clk}\) donne les quatre identités. Pour deux routes positives, le rapport des fréquences égale le rapport des masses ; la propriété multiplicative du caractère conserve la signature et les tours relatives avec \(R_{\rm act}^{\kappa_Y^r-\kappa_X^r}\).}

\end{proof}

\subsection{\HMTLang{Rigidité conjointe de l’action et de la lecture radiale}{Rigidité conjointe de l’action et de la lecture radiale}
}
\HMTLang{Dans cette sous-section, nous abrégeons \(\hbar=\hbar_{\rm ret}\). Soit \(M=M^{\beta,r}\) l’opérateur de masse déjà construit sur une fibre massive finie et strictement positive. Nous définissons, après \(M\),}{Dans cette sous-section, nous abrégeons \(\hbar=\hbar_{\rm ret}\). Soit \(M=M^{\beta,r}\) l’opérateur de masse déjà construit sur une fibre massive finie et strictement positive. Après \(M\), définissons}

\[
 X=\frac{M}{m_{\rm clk}},\qquad
 H=Mc_{\rm int}^2=Q_LX,\qquad
 \Lambda=L_\alpha X^{-1},\qquad
 \mathcal R_g=L_\alpha X.
\]
\begin{theorem}[\HMTLang{Deux relations conservées et coefficient gravitationnel unique}{Deux relations conservées et coefficient gravitationnel unique}
]
\HMTLang{Les lectures précédentes satisfont}{Les lectures précédentes satisfont}
\begingroup

\renewcommand{\theHequation}{delta22.R01}
\begin{equation}
 H\Lambda=\hbar c_{\rm int}I.
 \tag{R01}\label{delta22:vi:R01}
\end{equation}
\endgroup
\begingroup
\renewcommand{\theHequation}{delta22.R02}
\begin{equation}
 \mathcal R_g\Lambda=\Lambda\mathcal R_g=L_\alpha^2I.
 \tag{R02}\label{delta22:vi:R02}
\end{equation}
\endgroup
\HMTLang{Si la réalisation physique identifie ultérieurement la lecture radiale réduite au moyen de \(\mathcal R_g=(G/c_{\rm int}^2)M\), alors il existe un unique coefficient compatible :}{Si la réalisation physique identifie ultérieurement la lecture radiale réduite par \(\mathcal R_g=(G/c_{\rm int}^2)M\), alors il existe un unique coefficient compatible :}
\begingroup

\renewcommand{\theHequation}{delta22.G.unica}
\begin{equation}
 \boxed{G=\frac{c_{\rm int}^3L_\alpha^2}{\hbar}.}
 \label{delta22:vi:G-unica}
\end{equation}
\endgroup
\end{theorem}
\begin{proof}
\HMTLang{Les définitions et \(Q_L=\hbar c_{\rm int}/L_\alpha\) donnent directement (R01) et (R02). On en obtient \(\mathcal R_g=L_\alpha^2H/(\hbar c_{\rm int})\). Substituer \(H=Mc_{\rm int}^2\) et le comparer à \(\mathcal R_g=(G/c_{\rm int}^2)M\) produit \eqref{delta22:vi:G-unica}. Si \(G'\) satisfait les mêmes lectures, \((G'-G)M/c_{\rm int}^2=0\) ; l’inversibilité de \(M\) impose \(G'=G\). Le secteur nul est traité séparément et n’est pas inversé. En particulier, \(G\) se déduit après l’opérateur de masse et n’intervient pas dans sa construction.}{Les définitions et \(Q_L=\hbar c_{\rm int}/L_\alpha\) donnent directement (R01) et (R02). Elles impliquent \(\mathcal R_g=L_\alpha^2H/(\hbar c_{\rm int})\). La substitution de \(H=Mc_{\rm int}^2\) et la comparaison avec \(\mathcal R_g=(G/c_{\rm int}^2)M\) donnent \eqref{delta22:vi:G-unica}. Si \(G'\) satisfait les mêmes lectures, alors \((G'-G)M/c_{\rm int}^2=0\) ; l’inversibilité de \(M\) impose \(G'=G\). Le secteur nul est traité séparément et n’est pas inversé. En particulier, \(G\) est déduit après l’opérateur de masse et n’intervient pas dans sa construction.}

\end{proof}

\HMTLang{Si \(X\) est seulement semi-défini, on se restreint à \((\ker X)^\perp\). Avec la pseudo-inverse \(X^+\), \(\Lambda^+=L_\alpha X^+\) et \(P=P_{(\ker X)^\perp}\), les relations deviennent \(H\Lambda^+=\hbar c_{\rm int}P\) et \(\mathcal R_g\Lambda^+=L_\alpha^2P\). Le secteur nul conserve son traitement propre et ne produit pas de détermination vide de \(G\).}{Si \(X\) est seulement semi-défini positif, la construction est restreinte à \((\ker X)^\perp\). Avec la pseudo-inverse \(X^+\), \(\Lambda^+=L_\alpha X^+\) et \(P=P_{(\ker X)^\perp}\), les relations deviennent \(H\Lambda^+=\hbar c_{\rm int}P\) et \(\mathcal R_g\Lambda^+=L_\alpha^2P\). Le secteur nul conserve son traitement séparé et ne donne pas de détermination vide de \(G\).}


\HMTLang{Pour toute section d’action positive \(\hbar_a\), soit \(Q_{L,a}:=\hbar_a c_{\rm int}/L_\alpha=E_{P,a}\) l’énergie longitudinale de Planck dans cette même section. La même preuve donne \(G_a=c_{\rm int}^3L_\alpha^2/\hbar_a\) et conserve \(\hbar_aG_a/c_{\rm int}^3=L_\alpha^2\). La représentation au moyen de l’horloge est un corollaire, non la prémisse de la déduction :}{Pour toute section d’action positive \(\hbar_a\), soit \(Q_{L,a}:=\hbar_a c_{\rm int}/L_\alpha=E_{P,a}\) l’énergie longitudinale de Planck dans cette même section. La même preuve donne \(G_a=c_{\rm int}^3L_\alpha^2/\hbar_a\) et conserve \(\hbar_aG_a/c_{\rm int}^3=L_\alpha^2\). La représentation par l’horloge est un corollaire, non la prémisse de la déduction :}

\begin{corollary}[\HMTLang{Forme circulaire}{Forme circulaire}
]
\begin{equation}
 t_0=\frac{\pi_{\rm HMT}L_\alpha}{54c_{\rm int}}
 \quad\Longrightarrow\quad
 G_a=\left(\frac{54}{\pi_{\rm HMT}}\right)^2
 \frac{c_{\rm int}^5t_0^2}{\hbar_a}.
 \label{delta22:vi:G-circular}
\end{equation}
\end{corollary}

\begin{theorem}[\HMTLang{Réciprocité des lectures radiales}{Réciprocité des lectures radiales}
]
\HMTLang{Avec la convention réduite \(r_{g,\Gamma}^{(a)}=G_am_\Gamma/c_{\rm int}^2\),}{Avec la convention réduite \(r_{g,\Gamma}^{(a)}=G_am_\Gamma/c_{\rm int}^2\),}

\[
 r_{g,\Gamma}^{(a)}=\frac{\hbar_{\rm ret}}{\hbar_a}R_{108}\Xi_\Gamma,
 \qquad
 r_{g,\Gamma}^{(a)}\bar\lambda_\Gamma
 =\frac{\hbar_{\rm ret}}{\hbar_a}R_{108}^2.
\]
\HMTLang{Dans la même section \(a={\rm ret}\),}{Dans la même section \(a={\rm ret}\),}

\begin{equation}
 r_{g,\Gamma}=R_{108}\Xi_\Gamma,\qquad
 \bar\lambda_\Gamma=R_{108}/\Xi_\Gamma,\qquad
 r_{g,\Gamma}\bar\lambda_\Gamma=R_{108}^2.
 \label{delta22:vi:radios}
\end{equation}
\end{theorem}
\begin{proof}
\HMTLang{Substituer \(G_a\) et \(m_\Gamma=\hbar_{\rm ret}\omega_0\Xi_\Gamma/c_{\rm int}^2\) donne \(r_g^{(a)}=(\hbar_{\rm ret}/\hbar_a)R_{108}^2\omega_0\Xi_\Gamma/c_{\rm int}\). L’identité \(R_{108}\omega_0=c_{\rm int}\) prouve la première formule ; multiplier par \(R_{108}/\Xi_\Gamma\) donne la seconde. Le choix d’une même section annule le quotient des actions.}{Substituer \(G_a\) et \(m_\Gamma=\hbar_{\rm ret}\omega_0\Xi_\Gamma/c_{\rm int}^2\) donne \(r_g^{(a)}=(\hbar_{\rm ret}/\hbar_a)R_{108}^2\omega_0\Xi_\Gamma/c_{\rm int}\). L’identité \(R_{108}\omega_0=c_{\rm int}\) prouve la première formule ; la multiplication par \(R_{108}/\Xi_\Gamma\) donne la seconde. Choisir la même section annule le rapport des actions.}

\end{proof}
\HMTLang{La convention \(2Gm/c^2\) introduirait un facteur deux dans le produit et constituerait un autre lecteur. Ici, les longueurs sont des lectures du même état : dans les coordonnées logarithmiques, leurs déplacements sont \(\log\Xi_\Gamma\) et \(-\log\Xi_\Gamma\). Leur moyenne géométrique est \(R_{108}=L_\alpha\). L’identification physique de \(\mathcal R_g\) est expressément déclarée dans le théorème précédent ; elle ne redéfinit pas l’opérateur massique.}{La convention \(2Gm/c^2\) introduirait un facteur deux dans le produit et constituerait un lecteur différent. Ici, les longueurs sont des lectures du même état : dans les coordonnées logarithmiques, leurs déplacements sont \(\log\Xi_\Gamma\) et \(-\log\Xi_\Gamma\). Leur moyenne géométrique est \(R_{108}=L_\alpha\). L’identification physique de \(\mathcal R_g\) est explicitement énoncée dans le théorème précédent ; elle ne redéfinit pas l’opérateur de masse.}


\subsection{\HMTLang{Lecture thermique du spectre et compatibilité conjointe}{Lecture thermique du spectre et compatibilité conjointe}
}
\HMTLang{La transduction thermique est une application linéaire positive entre les droites dimensionnelles orientées de température et d’énergie : \(k_B:\mathcal L_\Theta\to\mathcal L_E\). La distribution uniforme des trois états locaux a pour information \(-3(1/3)\log(1/3)=\log3\). Une section positive de température \(\Theta_{108}\), normalisée par l’énergie d’une décision chronologique, détermine de manière unique le transducteur au moyen de}{La transduction thermique est une application linéaire positive entre les droites dimensionnelles orientées de température et d’énergie : \(k_B:\mathcal L_\Theta\to\mathcal L_E\). La distribution uniforme sur les trois états locaux a pour information \(-3(1/3)\log(1/3)=\log3\). Une section positive de température \(\Theta_{108}\), normalisée par l’énergie d’une décision chronologique, détermine de manière unique le transducteur au moyen de}

\[
 k_B\Theta_{108}\log3=E_{\rm clk}.
\]
\HMTLang{Dans des bases positives \(u_E,u_\Theta\), soient \(E_{\rm clk}=\epsilon u_E\), \(\Theta_{108}=\theta u_\Theta\) et \(k_B(u_\Theta)=b u_E\). L’équation donne \(b=\epsilon/(\theta\log3)>0\), ce qui prouve l’existence et l’unicité. Sous \(u'_E=s_Eu_E\), \(u'_\Theta=s_\Theta u_\Theta\), on a \(b'=(s_\Theta/s_E)b\), \(\epsilon'=\epsilon/s_E\), \(\theta'=\theta/s_\Theta\) ; l’identité est conservée. La coordonnée individuelle du transducteur maintient ces bases et la section thermique comme données déclarées. L’action et la période fixent préalablement \(E_{\rm clk}=h_{\rm ret}/(108t_0)\).}{Dans des bases positives \(u_E,u_\Theta\), soient \(E_{\rm clk}=\epsilon u_E\), \(\Theta_{108}=\theta u_\Theta\) et \(k_B(u_\Theta)=b u_E\). L’équation donne \(b=\epsilon/(\theta\log3)>0\), prouvant l’existence et l’unicité. Sous \(u'_E=s_Eu_E\), \(u'_\Theta=s_\Theta u_\Theta\), on a \(b'=(s_\Theta/s_E)b\), \(\epsilon'=\epsilon/s_E\) et \(\theta'=\theta/s_\Theta\) ; l’identité est conservée. La coordonnée individuelle du transducteur conserve ces bases et la section thermique comme données déclarées. L’action et la période fixent préalablement \(E_{\rm clk}=h_{\rm ret}/(108t_0)\).}
\HMTLang{Dans la réalisation canonique à cette température, avec une origine énergétique commune et un potentiel chimique nul,}{Dans la réalisation canonique à cette température, avec une origine énergétique commune et un potentiel chimique nul,}

\[
 \beta_{108}=(k_B\Theta_{108})^{-1},\qquad
 \beta_{108}E_\Gamma=(\log3)\Xi_\Gamma,\qquad
 e^{-\beta_{108}E_\Gamma}=3^{-\Xi_\Gamma}.
\]
\HMTLang{La preuve substitue \(E_\Gamma=E_{\rm clk}\Xi_\Gamma\). Pour la fibre finie, son expression opératorielle complète est}{La preuve substitue \(E_\Gamma=E_{\rm clk}\Xi_\Gamma\). Pour la fibre finie, son expression opératorielle complète est}

\[
 S(\rho):=-\operatorname{Tr}(\rho\log\rho),\qquad
 Z=\operatorname{Tr}(3^{-X}),\qquad
 \rho_{108}=3^{-X}/Z.
\]
\[
 S(\rho_{108})=\log Z+(\log3)\operatorname{Tr}(\rho_{108}X).
\]
\HMTLang{La dernière égalité provient de \(\log\rho_{108}=-(\log3)X-(\log Z)I\) et de \(\operatorname{Tr}\rho_{108}=1\). Si les valeurs propres sont regroupées avec des multiplicités physiques \(g_\Gamma\), \(Z=\sum_\Gamma g_\Gamma3^{-\Xi_\Gamma}\). On compte les états de la réalisation ; les descriptions équivalentes d’une histoire conservent leur identification. En dimension infinie, on exige \(3^{-X}\) de classe trace et, pour une entropie finie, \(\operatorname{Tr}(X3^{-X})<\infty\).}{La dernière égalité découle de \(\log\rho_{108}=-(\log3)X-(\log Z)I\) et de \(\operatorname{Tr}\rho_{108}=1\). Si les valeurs propres sont regroupées avec des multiplicités physiques \(g_\Gamma\), alors \(Z=\sum_\Gamma g_\Gamma3^{-\Xi_\Gamma}\). Les états de la réalisation sont comptés ; les descriptions équivalentes d’une histoire conservent leur identification. En dimension infinie, \(3^{-X}\) doit être de classe trace et une entropie finie requiert en outre \(\operatorname{Tr}(X3^{-X})<\infty\).}


\HMTLang{Ici, \(S\) est l’entropie adimensionnelle de von Neumann ; l’entropie dimensionnelle s’obtient au moyen du transducteur \(k_B\).}{Ici, \(S\) est l’entropie adimensionnelle de von Neumann ; l’entropie dimensionnelle s’obtient au moyen du transducteur \(k_B\).}


\HMTLang{La composition complète, dans une même section d’action et en utilisant \(R_{108}=L_\alpha\), devient}{La composition complète, sur la même section d’action et en utilisant \(R_{108}=L_\alpha\), est}

\begin{equation}
 \boxed{\Xi_\Gamma=\frac{m_\Gamma}{m_{\rm clk}}
 =\frac{\Omega_\Gamma}{\omega_0}
 =\frac{r_{g,\Gamma}}{L_\alpha}
 =\frac{L_\alpha}{\bar\lambda_\Gamma}
 =\frac{E_\Gamma}{k_B\Theta_{108}\log3}.}
 \label{delta22:vi:compatibilidad}
\end{equation}
\HMTLang{Cette égalité identifie les lectures d’une même section ; la généalogie qui la produit demeure dans \(b_e\), le caractère raffiné et l’exposant de retour. En \(\Xi=1\), les lectures radiales coïncident et l’énergie est \(E_{\rm clk}\) ; l’existence d’une espèce en ce point dépend du spectre effectivement construit. L’identité conjointe fournit un contrôle de compatibilité entre les réalisations dimensionnelles et conserve la distinction entre énergie propre, action, température et mémoire de route.}{Cette égalité identifie les lectures de la même section ; la généalogie qui la produit demeure dans \(b_e\), le caractère raffiné et l’exposant de retour. En \(\Xi=1\), les lectures radiales coïncident et l’énergie est \(E_{\rm clk}\) ; la présence d’une espèce en ce point dépend du spectre effectivement construit. L’identité conjointe fournit un contrôle de compatibilité entre les réalisations dimensionnelles et conserve la distinction entre énergie propre, action, température et mémoire de route.}

